% 2022 新高考Ⅰ卷 19：直三棱柱 ABC-A1B1C1（D 为 A1C 的中点）
\begin{tikzpicture}[scale=0.95, font=\small]
  \coordinate (A)  at (0,0.5);
  \coordinate (B)  at (1.15,-0.42);
  \coordinate (C)  at (2.5,0.5);
  \coordinate (A1) at (0,2.7);
  \coordinate (B1) at (1.15,1.78);
  \coordinate (C1) at (2.5,2.7);
  \coordinate (D)  at ($(A1)!0.5!(C)$);
  \draw (A) -- (B) -- (C);
  \draw[dashed] (A) -- (C);
  \draw (A) -- (A1)  (B) -- (B1)  (C) -- (C1);
  \draw (A1) -- (B1) -- (C1) -- cycle;
  \draw (A1) -- (B);
  \draw[dashed] (A1) -- (C);
  \draw[dashed] (B) -- (D);
  \fill (D) circle (0.9pt);
  \node[left] at (A) {$A$};
  \node[below] at (B) {$B$};
  \node[right] at (C) {$C$};
  \node[above left, inner sep=1pt] at (A1) {$A_1$};
  \node[below, inner sep=2.5pt] at (B1) {$B_1$};
  \node[above right, inner sep=1pt] at (C1) {$C_1$};
  \node[above right, inner sep=0.5pt] at (D) {$D$};
\end{tikzpicture}
